\subsection{DEFINITION OF THE EXTERIOR ALGEBRA OF A MODULE}

DEFINITION 1. Let A be a commutative ring and M an A-module. The exterior algebra of M, denoted by \(\bigwedge (\mathbf{M})\) or \(\operatorname {Alt}(\mathbf{M})\) or \(\bigwedge_{\mathbf{A}}(\mathbf{M})\) , is the algebra over A the quotient of the tensor algebra \(\mathsf{T}(\mathbf{M})\) by the two-sided ideal \(\mathfrak{J}''\) (also denoted by \(\mathfrak{J}_{\mathbf{M}}''\) ) generated by the elements \(x\otimes x\) , where \(x\) runs through M.

Since the ideal \(\mathfrak{J}''\) is generated by homogeneous elements of degree 2, it is a graded ideal (II, §11, no. 3, Proposition 2); we write \(\mathfrak{J}_n'' = \mathfrak{J}'' \cap \mathsf{T}^n(\mathbf{M})\) ; the algebra \(\bigwedge(\mathbf{M})\) is therefore graded by the graduation (called canonical) consisting of the \(\bigwedge^n(\mathbf{M}) = \mathsf{T}^n(\mathbf{M}) / \mathfrak{J}_n''\) . Then \(\mathfrak{J}_0'' = \mathfrak{J}_1'' = \{0\}\) and hence \(\bigwedge^0(\mathbf{M})\) is identified with A and \(\bigwedge^1(\mathbf{M})\) with \(\mathsf{T}^1(\mathbf{M}) = \mathbf{M}\) ; in what follows we shall always make these identifications and the canonical injection \(\mathbf{M} \to \bigwedge(\mathbf{M})\) will be denoted by \(\phi''\) or \(\phi_{\mathbf{M}}''\) .

PROPOSITION 1. Let E be an A-algebra and \(f: \mathbf{M} \to \mathbf{E}\) an A-linear mapping such that

\[
(f (x)) ^ {2} = 0 \quad \text { for all } x \in \mathbf {M}.\tag{1}
\]

There exists one and only one A-algebra homomorphism \(g: \bigwedge(\mathbf{M}) \to \mathbf{E}\) such that \(f = g \circ \phi''\) .

In other words, \((\bigwedge (\mathbf{M}),\phi '')\) is a solution of the universal mapping problem (Set Theory, IV, § 3, no. 1), where \(\Sigma\) is the species of A-algebra structure, the \(\alpha\) -mappings being the linear mappings from the A-module M to an A-algebra satisfying (1).

The uniqueness of \(g\) follows from the fact that \(\phi''(\mathbf{M}) = \mathbf{M}\) generates \(\bigwedge (\mathbf{M})\) . To prove the existence of \(g\) , we note that by § 5, no. 1, Proposition 1 there exists an A-algebra homomorphism \(g_1: \mathsf{T}(\mathbf{M}) \to \mathbf{E}\) such that \(f = g_1 \circ \phi\) ; we need to prove that \(g_1\) is zero on the ideal \(\mathfrak{J}''\) , for then if

\[
p \colon \mathbf {T} (\mathbf {M}) \rightarrow \bigwedge (\mathbf {M}) = \mathbf {T} (\mathbf {M}) / \mathfrak {J} ^ {\prime \prime}
\]

is the canonical homomorphism, we can write \(g_1 = g \circ p\) , where \(g: \bigwedge (\mathbf{M}) \to \mathbf{E}\) is an algebra homomorphism and the conclusion will follow from the fact that \(p \circ \phi = \phi''\) . Now, the kernel of \(g_1\) is a two-sided ideal which, by virtue of (1) and the relation \(g_1 \circ \phi = f\) , contains the elements \(x \otimes x\) for \(x \in \mathbf{M}\) . This completes the proof.

Remarks. (1) Suppose that E is a graded A-algebra of type Z, of graduation \((\mathbf{E}_{n})\) , and suppose also that the linear mapping f (assumed to satisfy (1)) is such that

\[
f (\mathbf {M}) \subset \mathbf {E} _ {1}.\tag{2}
\]

Then the relation \(g(x_1x_2 \ldots x_p) = f(x_1)f(x_2) \ldots f(x_p)\) with the \(x_i \in \mathbf{M}\) shows that \(g(\bigwedge^p (\mathbf{M})) \subset \mathrm{E}_p\) for all \(p \geqslant 0\) and hence \(g\) is a graded algebra homomorphism.

\begin{enumerate}
\def\labelenumi{(\arabic{enumi})}
\setcounter{enumi}{1}
\tightlist
\item
  To avoid confusion, the product of two elements u, v of the exterior algebra \(\bigwedge(\mathbf{M})\) is usually denoted by \(u \wedge v\) and is called the exterior product of u by v. The elements of \(\bigwedge^{n}(\mathbf{M})\) are therefore the sums of elements of the form \(x_{1} \wedge x_{2} \wedge \cdots \wedge x_{n}\) with \(x_{i} \in M\) for \(1 \leqslant i \leqslant n\) and are often called n-vectors.
\end{enumerate}

